% =============================================================================
\section{Side by side: PL7 and Kala Chakra 1.5}
\label{sec:compare}
% =============================================================================

Kala Chakra~1.5 here means the single-file app \texttt{Kala Chakra v1.5.html} (and its
sibling \texttt{KAAST Caster v1.5.html}) in \texttt{Advanced\_Setup/web/dist/}. Its casting
engine is \texttt{Advanced\_Setup/web/src/ephemeris.js}, \texttt{kp\_logic.js} and
\texttt{dasha.js}; I checked that the dist files contain these sources byte for byte.
\autoref{tab:compare} sets the two engines side by side.

\begin{table}[htbp]
  \centering
  \footnotesize
  \caption[PL7 versus Kala Chakra 1.5]{The two casting engines compared. Grades refer to the
  PL7 column; the Kala Chakra column is from its source code.}
  \label{tab:compare}
  \begin{tabularx}{\textwidth}{>{\raggedright\arraybackslash}p{2.5cm} >{\raggedright\arraybackslash}X >{\raggedright\arraybackslash}X}
    \toprule
    Item & Parashara's Light 7.0 & Kala Chakra 1.5 \\
    \midrule
    Ephemeris & Own files (\texttt{ephemdat.bin}, \texttt{topos*.dat}), J2000 equatorial; precise to end of 2048, coarse fallback after \tagNUM & Swiss Ephemeris 2.10.03 (WASM), files 1800--2399 \\
    $\Delta T$ & 1993 table, extrapolated: 83~s (2020) to 109~s (2048) \tagNUM & Solar Fire anchors 76--88~s (\emph{sf}) or Swiss $\approx 69$~s (\emph{modern}) \\
    Ayanamsa & 8 options; default Lahiri $=$ Swiss Lahiri $\pm 0.3''$; two KP variants \tagBIN & Lahiri only, Swiss $- 1.251''$ (Solar Fire clone) \\
    Equinox & Mean equinox of date \tagNUM & Mean equinox of date (Swiss sidereal mode) \\
    Node & True (in this installation; true/mean option exists) \tagNUM & True \\
    Planet positions & Geocentric; in this copy displaced by $\approx 0.97\times$ the Moon vector (up to $\sim\!10'$) \tagNUM & Geocentric apparent, $\le 1''$ from Solar Fire \\
    House cusps & Placidus, geographic latitude \tagBIN/\tagINF & Topocentric, geocentric (WGS84) latitude \\
    Bhava chart & Sripati (centres, quadrant trisection) and equal \tagDOC & Not offered (houses are KP beginnings) \\
    Rounding & Not determined & Truncated to whole arcseconds \\
    Star/sub lords & Star and sub \tagDOC & Star plus five-level sub chain \\
    KP extras & Significators, ruling planets, Fortuna, horary 1--249 \tagDOC/\tagINF & Cusp-number vectors, span rule, power calculation \\
    Vargas & 16 Parashari $+$ 4 Tajika $+$ D150 $+$ 10 alternates \tagBIN & None \\
    Yogi points & Yogi $=$ Sun $+$ Moon $+ 93^\circ 20'$ \tagBIN & None \\
    Dashas & $\ge 20$ systems, 5 levels; ``365.25'' option actually 365.2421875~d \tagBIN & Vimshottari, 6 levels, 365.25~d \\
    Jaimini \& padas & Karakas (7/8), rashi drishti, arudhas, upapada, special lagnas, argala \tagDOC & None \\
    Transit viewer & Animated circle, steps 1~s to 10~yr, via a 12-hourly month cache \tagBIN/\tagNUM & None (one chart per cast) \\
    Guna matching & Ashtakoota, Dashakoota, Manglik \tagDOC & None \\
    \bottomrule
  \end{tabularx}
\end{table}

\subsection{What the differences do to star and sub lords}

A KP sub is at least $40'$ wide (the Sun's sub, $13^\circ 20' \times 6/120$), so a difference
matters only for points lying within that difference of a boundary. Using random longitudes and the cusp
simulation of \autoref{tab:cusps} \tagNUM:

\begin{itemize}[leftmargin=*]
  \item PL7 Lahiri vs Kala Chakra Lahiri ($1.4''$): 0.02\,\% of points change star or sub
  lord.
  \item PL7 ``Krishnamurti'' vs Kala Chakra ($27''$): 0.5\,\%.
  \item PL7 ``Krishnamurti New'' vs Kala Chakra ($5'59''$): 6.8\,\%.
  \item Placidus with geographic latitude vs Topocentric with geocentric latitude: 0.1\,\%
  of cusps at Singapore and 1.7\,\% at Lapanga.
  \item The Moon-vector error ($\sim\!1'$--$10'$) in this PL7 copy: for an error of
  about $8'$ against an average sub width of $89'$, roughly one Sun, Mercury or Venus position
  in ten gets a different sub lord; the Moon and the nodes are unaffected.
\end{itemize}

\begin{keybox}\kidea For the family's KP horary work the casting engines differ in three
choices, Lahiri versus KP ayanamsa, Placidus versus Topocentric, and the latitude
convention, plus one defect specific to this PL7 copy. Kala Chakra's choices were made
deliberately to replicate the Solar Fire profile the family uses, and they are validated
against the family's own hand-calculated charts. That validation, not either program, is
the authority for the method.\end{keybox}

% =============================================================================
\section{Tests you can run in PL7 yourself}
\label{sec:tests}
% =============================================================================

These take a few minutes each and settle the points this report could only infer.

\begin{enumerate}[leftmargin=*]
  \item \textbf{Is this copy's Sun wrong?} Cast one chart in PL7 with Lahiri, say for
  1~January~2020 at 0:00~UT (5:30~IST) at any place, and read the Sun to the second.
  Cast the same moment in KAAST Caster in \emph{modern} mode, which gives about
  $255^\circ 52' 40''$ (Sagittarius $15^\circ 52' 40''$). PL7's own transit cache holds
  $256^\circ 01' 09''$ for that instant, $8'29''$ larger. If PL7's chart shows the larger value,
  the Moon-vector error is present in natal charts too. The Moon should agree to about $10''$
  (PL7's larger $\Delta T$ accounts for about $7''$).
  \item \textbf{Which KP ayanamsa does the KP printout use?} Set the Calculation tab to
  Lahiri, print the Krishnamurti chart, and read the ``KP Ayanamsha'' line. Compare with
  \autoref{tab:ayanamsa}: about $24^\circ 11' 54''$ (Krishnamurti) or $24^\circ 06' 23''$
  (Krishnamurti New) at 1~January~2025.
  \item \textbf{Rahu true or mean.} In Options $\to$ Calculation $\to$ Miscellaneous, check
  the ``Rahu \& Ketu'' item; the cache says \emph{true} in this installation.
  \item \textbf{D11.} Display the Ekadashamsha and its alternate (II/III) version for a
  planet at Taurus $1^\circ$: the main D11 should put it in Pisces, the alternate in Aries.
\end{enumerate}
